\section{Robustness to model error}\label{sec:robustness}

The bounds assume that the fitting model is exact. \Cref{tab:robustness} generates noise-free
data with an effect the ideal model ignores and fits them with the ideal two-scan model (complex,
all seven parameters free, ten starting points, $\Bone$ restricted to the fitting range $[0.6,1.4]$) for
WM, GM and CSF at $\Bone=0.7,0.8,\dots,1.3$ and at the transmit scales where the design's scan-2 angle is
$330^\circ$ and $345^\circ$.
\begin{itemize}
\item \emph{Pulse ratio}: the scan-2 angle $1\%$ larger than nominal relative to scan 1, as produced
  by an imperfect calibration of two pulses of different shape or duration.
\item \emph{Slab-edge ramp}: the flip angle falls linearly from $100\%$ to $90\%$ across the voxel, each
  position in its own steady state, as at the edge of a 3D slab.
\item \emph{Finite pulses}: hard pulses at a fixed peak $\Bone$, so that the duration is proportional to
  the nominal angle ($1$\,ms for $330^\circ$), with relaxation during the pulse; echo times are measured
  from the pulse centre. On resonance, and with $50$\,Hz off resonance (precession during the pulse).
\end{itemize}
Besides the bias we give the probability that the misfit would be flagged at SNR 300 by a $\chi^2$
test on the residual (level $99.9\%$, 29 degrees of freedom; noncentral $\chi^2$ with the noise-free
residual as noncentrality).

\begin{table}[!htbp]
\centering
\caption{Bias of $\Tone$, $100(\hat\Tone/\Tone-1)$ (\%), when the ideal two-scan model is fitted to data with
the effect (noise-free, all parameters free). Columns 3--5: at $\Bone=1$; column 6: largest $|$bias$|$ over
WM, GM, CSF and the sampled $\Bone$ ($0.7,0.8,\dots,1.3$ and the scan-2 angles $330^\circ$, $345^\circ$), a lower bound
for the continuous range; column 7: largest $|\delta\ln\Tone|$ divided by the $\Tone$
standard deviation at SNR$(\Mzero)=300$ (bound/300); column 8: probability that the misfit of the
largest-bias case is flagged at SNR 300. $^*$: a fit ended on a safeguard bound of the optimiser.
Generated by \texttt{scripts/flip\_design\_robustness.py}.}
\label{tab:robustness}
\small\setlength{\tabcolsep}{4pt}
\begin{tabular}{llcccccc}
\toprule
& & \multicolumn{3}{c}{$\Bone=1$} & \multicolumn{3}{c}{largest over the sampled cases}\\
\cmidrule(lr){3-5}\cmidrule(lr){6-8}
$\alpha_1/\alpha_2$ & effect & WM & GM & CSF & $|$bias$|$ (\%) & $|$bias$|$/SD & flagged\\
\midrule
15/330 & pulse ratio $+1\%$ & $-1.8$ & $-1.8$ & $-1.8$ & 2.3 & 0.2 & 0.00\\
 & slab-edge ramp & $+1.7$ & $+0.9$ & $+0.7$ & 40.3 & 1.3 & 0.58\\
 & finite pulses & $+19.7$ & $+18.5$ & $+5.4$ & 19.7 & 2.0 & 0.04\\
 & finite pulses, 50\,Hz & $+22.2$ & $+19.6$ & $+11.9$ & 22.2 & 2.2 & 0.99\\
\midrule
20/276 & pulse ratio $+1\%$ & $-1.7$ & $-1.7$ & $-1.7$ & 2.1 & 0.2 & 0.00\\
 & slab-edge ramp & $-0.3$ & $-0.3$ & $-0.2$ & 16.3 & 1.1 & 0.10\\
 & finite pulses & $+2.4$ & $+2.2$ & $+0.6$ & 13.8 & 1.2 & 0.01\\
 & finite pulses, 50\,Hz & $+2.7$ & $+2.5$ & $+3.8$ & 13.3 & 1.1 & 0.92\\
\midrule
14/309 & pulse ratio $+1\%$ & $-1.8$ & $-1.8$ & $-1.8$ & 2.2 & 0.2 & 0.00\\
 & slab-edge ramp$^*$ & $-0.1$ & $-0.3$ & $0.0$ & 74.1 & 3.9 & 0.83\\
 & finite pulses & $+6.4$ & $+5.1$ & $+0.6$ & 19.5 & 2.0 & 0.03\\
 & finite pulses, 50\,Hz & $+6.9$ & $+5.6$ & $+8.4$ & 21.9 & 2.2 & 0.98\\
\midrule
21/270 & pulse ratio $+1\%$ & $-1.7$ & $-1.7$ & $-1.7$ & 2.0 & 0.2 & 0.00\\
 & slab-edge ramp & $-0.3$ & $-0.3$ & $-0.1$ & 13.0 & 0.8 & 0.02\\
 & finite pulses & $+2.2$ & $+2.1$ & $+1.8$ & 14.4 & 1.2 & 0.01\\
 & finite pulses, 50\,Hz & $+2.5$ & $+2.4$ & $+3.5$ & 13.9 & 1.2 & 0.88\\
\midrule
17/254 & pulse ratio $+1\%$ & $-1.6$ & $-1.6$ & $-1.6$ & 2.0 & 0.2 & 0.00\\
 & slab-edge ramp & $-0.2$ & $-0.2$ & $0.0$ & 0.9 & 0.1 & 0.00\\
 & finite pulses & $+2.0$ & $+1.9$ & $+2.1$ & 10.0 & 0.8 & 0.05\\
 & finite pulses, 50\,Hz & $+2.3$ & $+2.2$ & $+3.1$ & 11.1 & 0.9 & 0.88\\
\midrule
28/180 & pulse ratio $+1\%$ & $-2.0$ & $-2.0$ & $-2.0$ & 3.1 & 0.1 & 0.00\\
 & slab-edge ramp & $-0.2$ & $-0.2$ & $0.0$ & 0.4 & 0.0 & 0.00\\
 & finite pulses & $+2.0$ & $+2.0$ & $+1.8$ & 3.4 & 0.1 & 0.00\\
 & finite pulses, 50\,Hz & $+2.3$ & $+2.3$ & $+2.2$ & 4.0 & 0.1 & 0.00\\

\bottomrule
\end{tabular}
\end{table}

\paragraph{Pulse ratio: the same for every design, and invisible.} A $1\%$ ratio error biases $\Tone$ by
$-1.6$ to $-2.0\%$ at $\Bone=1$ in every design (up to $-3.1\%$ over the $\Bone$ range), as \cref{prop:ratio}
predicts, and the fit absorbs it with zero residual. No choice of angles avoids it
($|\mathrm d\ln\Tone/\mathrm d\varepsilon|\ge1.64$): the ratio must be calibrated, to about $0.5\%$ if $\Tone$ is to be accurate to
$1\%$. The bias is only $0.2$ standard deviations per voxel at SNR 300, but it is systematic and does
not average out over a region.

\paragraph{Everything goes wrong near $360^\circ$.} The other two effects share one pattern: they bias
$\Tone$ strongly where the actual scan-2 angle approaches $360^\circ$, i.e.\ where its effective angle is small
(finite pulses most at $330$--$350^\circ$, flip-angle spreads at and across $360^\circ$), and hardly at all elsewhere; crossings of $180^\circ$ are harmless, because the scan-2 signal there is
small in any case ($|v|\approx\zeta^2|\delta|$ at $180^\circ+\delta$ against $|\delta|$ at $360^\circ+\delta$).
\begin{itemize}
\item \emph{Finite pulses.} With a $1$\,ms $330^\circ$ pulse the published design overestimates $\Tone$ by $20\%$
  in WM and $18.5\%$ in GM at $\Bone=1$, two standard deviations at SNR 300, with a misfit that is flagged
  in only $2$--$4\%$ of voxels. The margin design $21^\circ/270^\circ$ gives $2.2\%$ and $2.1\%$ there. The bias grows with
  the pulse duration: for a $330^\circ$ pulse of $0.125$, $0.25$, $0.5$, $1$, $2$ and $4$\,ms it is $3.4$, $6.6$, $12.7$,
  $19.7$, $22.6$ and $33.6\%$ for the published design (steep, then irregular) and $0.27$, $0.54$, $1.1$, $2.2$,
  $4.5$ and $9.6\%$ for $21^\circ/270^\circ$ (nearly proportional; durations scale with the nominal angle, so the
  $270^\circ$ pulse lasts $270/330$ of these). The advantage is local, though: every design meets the same
  problem where its scan-2 pulse passes $330$--$350^\circ$ (the margin design at $\Bone=1.28$, $-14\%$; $20^\circ/276^\circ$ at
  $\Bone=1.25$, $-14\%$; $17^\circ/254^\circ$ at $\Bone=1.3$, $+10\%$), so over the whole range the largest bias falls only from
  $20\%$ to $10$--$14\%$ for the no-crossing designs ($14^\circ/309^\circ$ reaches $19.5\%$ at its $330^\circ$ point; $3\%$ for the
  $180^\circ$ design, whose scan-2 pulse never gets there). Off resonance the
  misfit of the designs with $\alpha_2\ge254^\circ$ becomes detectable (flagged in $88$--$99\%$ of voxels in the
  worst case); the bias changes by at most $2.5$ points in WM and GM but grows in CSF (e.g.\ from
  $0.6$ to $8.4\%$ for $14^\circ/309^\circ$ at $\Bone=1$). Finite-pulse effects are deterministic and can be built into the
  forward model, as is done for balanced SSFP~\cite{bieri2009}; any design whose scan-2 pulse reaches
  $330$--$360^\circ$ over the $\Bone$ range needs it.
\item \emph{Slab-edge spread.} All designs estimate $\Bone$ as the voxel average ($\approx-5\%$) and $\Tone$ within
  $2\%$ as long as the scan-2 angle stays below about $330^\circ$. Near $360^\circ$ the bias grows: $-13\%$ and $-16\%$
  for $21^\circ/270^\circ$ and $20^\circ/276^\circ$ at $\Bone=1.3$ (scan-2 angle $351^\circ$ and $359^\circ$), and $+32$ to $+74\%$ where the
  spread straddles $360^\circ$ (the published design at $\Bone=1.1$--$1.2$, $14^\circ/309^\circ$ at $1.2$), because the two sides of
  the null contribute signals of opposite sign. These misfits are hard to detect: the $40\%$ CSF
  bias of the published design at $\Bone=1.2$ is flagged in $58\%$ of voxels, the $16\%$ bias of $20^\circ/276^\circ$ in
  $10\%$, and the margin design's $13\%$ in $2\%$. The two-stage design $17^\circ/254^\circ$, whose scan-2 angle never
  exceeds $330^\circ$, stays below $1\%$.
\end{itemize}

\paragraph{2D slice profiles.} An unmodelled 2D slice profile is a much larger error than the
slab-edge spread and is not tabulated: the scaled small-tip profile used for the slab edge is not
valid for a $250$--$330^\circ$ sinc pulse, whose rotation axis and tip angle vary across the slice in a
way only a Bloch simulation of the pulse captures. MPME relaxometry needs a 3D acquisition, as
in~\cite{cheng2019}, or the slice profile in the forward model.
